\begin{bt}%[9D6H18-2]
% nguon: Đề thi chính thức Tuyển sinh 10 TP.HCM năm học 2025-2026 – Bài 1
Cho hàm số $y = \dfrac{x^2}{2}$ có đồ thị $(P)$.
\begin{enumerate}[a)]
	\item Vẽ đồ thị $(P)$ trên hệ trục toạ độ.
	\item Tìm toạ độ các điểm thuộc $(P)$ có tung độ bằng $18$.
\end{enumerate}
\loigiai{
\begin{enumerate}[a)]
	\item Bảng giá trị tương ứng giữa $x$ và $y$ của hàm số $y = \dfrac{x^2}{2}$:
	\begin{center}
	\begin{tabular}{|c|c|c|c|c|c|}
	\hline
	$x$ & $-4$ & $-2$ & $0$ & $2$ & $4$ \\
	\hline
	$y = \dfrac{x^2}{2}$ & $8$ & $2$ & $0$ & $2$ & $8$ \\
	\hline
	\end{tabular}
	\end{center}
	Đồ thị $(P)$ là một parabol đi qua gốc toạ độ $O(0; 0)$, nhận trục tung $Oy$ làm trục đối xứng và nằm phía trên trục hoành.
	\item Điểm thuộc $(P)$ có tung độ bằng $18$ thoả mãn phương trình:
	$18 = \dfrac{x^2}{2} \Leftrightarrow x^2 = 36 \Leftrightarrow x = 6$ hoặc $x = -6$.\\
	Vậy toạ độ các điểm thoả mãn yêu cầu đề bài là $(6; 18)$ và $(-6; 18)$.
\end{enumerate}
}
\end{bt}

\begin{bt}%[9D6H20-2]
% nguon: Đề thi chính thức Tuyển sinh 10 TP.HCM năm học 2025-2026 – Bài 2
Cho phương trình $2x^2 - 7x + 4 = 0$.
\begin{enumerate}[a)]
	\item Chứng minh phương trình trên có hai nghiệm phân biệt $x_1, x_2$.
	\item Không giải phương trình, hãy tính giá trị của biểu thức $A = x_1(3x_2 + x_1) + x_2^2$.
\end{enumerate}
\loigiai{
\begin{enumerate}[a)]
	\item Xét phương trình $2x^2 - 7x + 4 = 0$, ta có:
	$\Delta = b^2 - 4ac = (-7)^2 - 4\cdot 2\cdot 4 = 49 - 32 = 17 > 0$.\\
	Vì $\Delta > 0$ nên phương trình luôn có hai nghiệm phân biệt $x_1, x_2$.
	\item Theo hệ thức Viète, ta có:
	$\begin{cases} x_1 + x_2 = -\dfrac{b}{a} = \dfrac{7}{2} \\ x_1 x_2 = \dfrac{c}{a} = \dfrac{4}{2} = 2. \end{cases}$\\
	Biến đổi biểu thức $A$:
	$A = x_1(3x_2 + x_1) + x_2^2 = 3x_1 x_2 + x_1^2 + x_2^2 = 3x_1 x_2 + (x_1 + x_2)^2 - 2x_1 x_2 = (x_1 + x_2)^2 + x_1 x_2$.\\
	Thay $x_1 + x_2 = \dfrac{7}{2}$ và $x_1 x_2 = 2$ vào biểu thức $A$, ta được:
	$A = \left(\dfrac{7}{2}\right)^2 + 2 = \dfrac{49}{4} + 2 = \dfrac{57}{4}$.\\
	Vậy $A = \dfrac{57}{4}$.
\end{enumerate}
}
\end{bt}

\begin{bt}%[9D8H26-3]
% nguon: Đề thi chính thức Tuyển sinh 10 TP.HCM năm học 2025-2026 – Bài 3
Biểu đồ tròn cho biết tỉ lệ về số lượng các loại bảo hiểm đã bán được trong tháng 4/2025 của một công ty. Biết rằng trong tháng này, công ty đã bán được $300$ gói bảo hiểm các loại cho $300$ khách hàng khác nhau:
\begin{itemize}
	\item Loại bảo hiểm A: $25\%$
	\item Loại bảo hiểm B: $15\%$
	\item Loại bảo hiểm C: $33\%$
	\item Loại bảo hiểm D: $27\%$
\end{itemize}
\begin{enumerate}[a)]
	\item Tính số lượng cụ thể của mỗi loại bảo hiểm mà công ty đã bán được trong tháng 4/2025?
	\item Bộ phận chăm sóc khách hàng chọn ngẫu nhiên một khách hàng đã mua bảo hiểm của công ty trong tháng 4/2025 để khảo sát. Tính xác suất của biến cố: ``Khách hàng được chọn không mua loại bảo hiểm B''.
\end{enumerate}
\loigiai{
\begin{enumerate}[a)]
	\item Số lượng cụ thể của từng loại bảo hiểm công ty đã bán được trong tháng 4/2025 là:
	\begin{itemize}
		\item Số lượng gói bảo hiểm loại A: $300 \cdot 25\% = 75$ (gói).
		\item Số lượng gói bảo hiểm loại B: $300 \cdot 15\% = 45$ (gói).
		\item Số lượng gói bảo hiểm loại C: $300 \cdot 33\% = 99$ (gói).
		\item Số lượng gói bảo hiểm loại D: $300 \cdot 27\% = 81$ (gói).
	\end{itemize}
	\item Số các kết quả có thể xảy ra khi chọn ngẫu nhiên 1 khách hàng là $n(\Omega) = 300$.\\
	Gọi biến cố $E$: ``Khách hàng được chọn không mua loại bảo hiểm B''.\\
	Khi đó khách hàng được chọn mua loại bảo hiểm A, C hoặc D. Số kết quả thuận lợi cho biến cố $E$ là:
	$n(E) = 75 + 99 + 81 = 255$ (khách hàng).\\
	Xác suất của biến cố $E$ là:
	$P(E) = \dfrac{n(E)}{n(\Omega)} = \dfrac{255}{300} = \dfrac{17}{20} = 0{,}85$.\\
	Vậy xác suất cần tìm là $\dfrac{17}{20}$.
\end{enumerate}
}
\end{bt}

\begin{bt}%[9D6V21-3]
% nguon: Đề thi chính thức Tuyển sinh 10 TP.HCM năm học 2025-2026 – Bài 4
Một khu vườn hình chữ nhật có chiều dài là $2x\text{ (m)}$ và chiều rộng là $x\text{ (m)}$ ($x > 4$). Bác Ba làm một lối đi xung quanh khu vườn rộng $2\text{ mét}$. Phần đất còn lại ở bên trong dùng để trồng hoa.
\begin{enumerate}[a)]
	\item Viết biểu thức theo $x$ biểu diễn diện tích phần đất dùng để trồng hoa và thu gọn biểu thức đó.
	\item Giả sử diện tích phần đất trồng hoa là $4800\text{ m}^2$. Tính chiều dài và chiều rộng của khu vườn.
\end{enumerate}
\loigiai{
\begin{enumerate}[a)]
	\item Vì lối đi quanh khu vườn rộng $2\text{ m}$ nên phần đất trồng hoa bên trong là hình chữ nhật có:
	\begin{itemize}
		\item Chiều rộng là: $x - 2\cdot 2 = x - 4\text{ (m)}$.
		\item Chiều dài là: $2x - 2\cdot 2 = 2x - 4\text{ (m)}$.
	\end{itemize}
	Biểu thức biểu diễn diện tích phần đất dùng để trồng hoa là:
	$S = (x - 4)(2x - 4) = 2x^2 - 4x - 8x + 16 = 2x^2 - 12x + 16\text{ (m}^2\text{)}$.
	\item Theo giả thiết, diện tích đất trồng hoa là $4800\text{ m}^2$ nên ta có phương trình:
	$2x^2 - 12x + 16 = 4800 \Leftrightarrow 2x^2 - 12x - 4784 = 0 \Leftrightarrow x^2 - 6x - 2392 = 0$.\\
	Ta có $\Delta' = (-3)^2 - 1\cdot(-2392) = 9 + 2392 = 2401 = 49^2 > 0$.\\
	Phương trình có hai nghiệm:
	$x_1 = 3 + 49 = 52$ (thoả mãn điều kiện $x > 4$),\\
	$x_2 = 3 - 49 = -46$ (loại vì $x > 4$).\\
	Với $x = 52$, ta có:
	\begin{itemize}
		\item Chiều rộng của khu vườn là: $x = 52\text{ (m)}$.
		\item Chiều dài của khu vườn là: $2x = 2\cdot 52 = 104\text{ (m)}$.
	\end{itemize}
	Vậy chiều dài của khu vườn là $104\text{ m}$ và chiều rộng là $52\text{ m}$.
\end{enumerate}
}
\end{bt}

\begin{bt}%[9H0V31-3]
% nguon: Đề thi chính thức Tuyển sinh 10 TP.HCM năm học 2025-2026 – Bài 5
Một hộp đựng bóng tennis có dạng hình trụ chứa vừa khít $4$ quả bóng tennis có dạng hình cầu. Biết diện tích bề mặt mỗi quả bóng tennis là $132{,}67\text{ cm}^2$.
Cho biết diện tích bề mặt hình cầu bán kính $R$ là $S = 4\pi R^2$. Diện tích toàn phần hình hộp chữ nhật có ba kích thước $a, b, c$ là $S_{tp} = 2(ab + bc + ca)$.
\begin{enumerate}[a)]
	\item Tính bán kính của mỗi quả bóng tennis.
	\item Nhà sản xuất sử dụng các thùng giấy hình hộp chữ nhật (có nắp) để chứa $12$ hộp tennis sao cho các hộp tennis được xếp vừa khít trong thùng giấy (xếp thành $3$ hàng, mỗi hàng $4$ hộp; hoặc $2$ hàng, mỗi hàng $6$ hộp theo chiều đứng). Hỏi cần tối thiểu bao nhiêu $\text{m}^2$ giấy để thiết kế một thùng như trên (giả sử mép dán không đáng kể). Các kết quả được làm tròn đến hàng phần trăm.
\end{enumerate}
\loigiai{
\begin{enumerate}[a)]
	\item Diện tích bề mặt mỗi quả bóng tennis là $132{,}67\text{ cm}^2$, do đó:
	$4\pi R^2 = 132{,}67 \Rightarrow R^2 = \dfrac{132{,}67}{4\pi} \approx \dfrac{132{,}67}{4\cdot 3{,}14159} \approx 10{,}557$.\\
	Suy ra bán kính của mỗi quả bóng tennis là:
	$R = \sqrt{\dfrac{132{,}67}{4\pi}} \approx 3{,}25\text{ (cm)}$.
	\item Đường kính của mỗi quả bóng tennis là:
	$d = 2R = 2\cdot 3{,}25 = 6{,}5\text{ (cm)}$.\\
	Vì hộp tennis hình trụ chứa vừa khít $4$ quả bóng nên:
	\begin{itemize}
		\item Đường kính đáy hộp hình trụ là: $D = d = 6{,}5\text{ (cm)}$.
		\item Chiều cao của hộp hình trụ là: $h = 4d = 4\cdot 6{,}5 = 26\text{ (cm)}$.
	\end{itemize}
	Để chứa vừa khít $12$ hộp tennis dựng đứng thành khối hình hộp chữ nhật, cách xếp tối ưu diện tích là xếp thành $3$ hàng, mỗi hàng $4$ hộp (hoặc đáy gồm $3 \times 4$ hộp):
	\begin{itemize}
		\item Chiều rộng thùng giấy là: $a = 3\cdot D = 3\cdot 6{,}5 = 19{,}5\text{ (cm)}$.
		\item Chiều dài thùng giấy là: $b = 4\cdot D = 4\cdot 6{,}5 = 26\text{ (cm)}$.
		\item Chiều cao thùng giấy là: $c = h = 26\text{ (cm)}$.
	\end{itemize}
	Diện tích toàn phần của thùng giấy hình hộp chữ nhật là:
	$S_{tp} = 2(ab + bc + ca) = 2(19{,}5\cdot 26 + 26\cdot 26 + 26\cdot 19{,}5) = 2(507 + 676 + 507) = 2\cdot 1690 = 3380\text{ (cm}^2\text{)}$.\\
	Đổi sang đơn vị mét vuông:
	$3380\text{ cm}^2 = \dfrac{3380}{10000}\text{ m}^2 = 0{,}338\text{ m}^2 \approx 0{,}34\text{ (m}^2\text{)}$.\\
	Vậy cần tối thiểu khoảng $0{,}34\text{ m}^2$ giấy để thiết kế một thùng như trên.
\end{enumerate}
}
\end{bt}

\begin{bt}%[9D1C3-1]
% nguon: Đề thi chính thức Tuyển sinh 10 TP.HCM năm học 2025-2026 – Bài 6
Từ vị trí $A$ của một công viên có dạng hình vuông $ABCD$ cạnh $a\text{ (km)}$, hai bạn Hoà và Bình bắt đầu chạy bộ cùng lúc với vận tốc không đổi dọc theo các cạnh của hình vuông và theo hai hướng khác nhau (Hoà chạy theo hướng $A \to B \to C \to D$, Bình chạy theo hướng $A \to D \to C \to B$). Biết rằng hai bạn gặp nhau lần thứ nhất tại vị trí $E$ trên cạnh $CD$ cách $A$ một khoảng bằng $1\text{ km}$ ($AE = 1\text{ km}$) và gặp lại nhau lần thứ hai tại vị trí $F$ trên cạnh $AB$ cách $A$ một khoảng bằng $0{,}4\text{ km}$ ($AF = 0{,}4\text{ km}$). Gọi $x, y\text{ (km/h)}$ lần lượt là vận tốc của Hoà và Bình.
\begin{enumerate}[a)]
	\item Chứng minh rằng $\dfrac{x}{y} = \dfrac{AB + BC + CE}{AD + DE}$.
	\item Tìm giá trị của $a$.
\end{enumerate}
\loigiai{
\begin{enumerate}[a)]
	\item Hai bạn xuất phát cùng lúc từ $A$. Khi gặp nhau lần thứ nhất tại $E$ trên cạnh $CD$:
	\begin{itemize}
		\item Quãng đường Hoà đã chạy là: $s_1 = AB + BC + CE$.
		\item Quãng đường Bình đã chạy là: $s_2 = AD + DE$.
	\end{itemize}
	Vì thời gian chạy của hai bạn từ lúc xuất phát đến lúc gặp nhau lần đầu là như nhau nên:
	$t_1 = \dfrac{s_1}{x} = \dfrac{s_2}{y} \Leftrightarrow \dfrac{AB + BC + CE}{x} = \dfrac{AD + DE}{y}$.\\
	Suy ra $\dfrac{x}{y} = \dfrac{AB + BC + CE}{AD + DE}$.
	\item Chu vi công viên hình vuông là $4a\text{ (km)}$.\\
	Tổng quãng đường hai bạn chạy được tính đến lần gặp thứ nhất là:
	$s_1 + s_2 = (AB + BC + CE) + (AD + DE) = AB + BC + CD + DA = 4a\text{ (km)}$.\\
	Do đó thời gian chạy đến lần gặp thứ nhất là:
	$t_1 = \dfrac{4a}{x + y}$.\\
	Từ lần gặp thứ nhất đến lần gặp thứ hai tại $F$, hai bạn tiếp tục chạy ngược chiều quanh chu vi hình vuông và gặp nhau, nên tổng quãng đường hai bạn chạy thêm đúng bằng một vòng chu vi hình vuông là $4a\text{ (km)}$.\\
	Do đó, thời gian từ lúc xuất phát đến khi gặp nhau lần thứ hai là:
	$t_2 = \dfrac{4a + 4a}{x + y} = \dfrac{8a}{x + y} = 2t_1$.\\
	Vì vận tốc của Bình là không đổi nên quãng đường Bình chạy đến lần gặp thứ hai gấp đôi quãng đường Bình chạy đến lần gặp thứ nhất.\\
	Tại lần gặp thứ hai ở $F$ trên cạnh $AB$ cách $A$ là $AF = 0{,}4\text{ km}$, quãng đường Bình đã chạy từ $A$ qua $D, C, B$ đến $F$ là:
	$s'_2 = AD + DC + CB + BF = 4a - AF = 4a - 0{,}4\text{ (km)}$.\\
	Do $s'_2 = 2s_2$, ta có:
	$4a - 0{,}4 = 2(AD + DE) = 2(a + DE) \Leftrightarrow 4a - 0{,}4 = 2a + 2DE \Leftrightarrow 2DE = 2a - 0{,}4 \Leftrightarrow DE = a - 0{,}2$.\\
	Tam giác $ADE$ vuông tại $D$ có $AD = a$, $DE = a - 0{,}2$ và cạnh huyền $AE = 1\text{ km}$. Áp dụng định lí Pythagore:
	$AD^2 + DE^2 = AE^2 \Leftrightarrow a^2 + (a - 0{,}2)^2 = 1^2 \Leftrightarrow a^2 + a^2 - 0{,}4a + 0{,}04 = 1 \Leftrightarrow 2a^2 - 0{,}4a - 0{,}96 = 0$.\\
	Chia cả hai vế cho $2$: $a^2 - 0{,}2a - 0{,}48 = 0$.\\
	Ta có $\Delta' = (-0{,}1)^2 - 1\cdot(-0{,}48) = 0{,}01 + 0{,}48 = 0{,}49 = (0{,}7)^2$.\\
	Phương trình có hai nghiệm:
	$a_1 = 0{,}1 + 0{,}7 = 0{,}8$ (thoả mãn $a > 0{,}2$),\\
	$a_2 = 0{,}1 - 0{,}7 = -0{,}6$ (loại vì $a > 0$).\\
	Vậy $a = 0{,}8\text{ (km)}$.
\end{enumerate}
}
\end{bt}

\begin{bt}%[9H9C29-3]
% nguon: Đề thi chính thức Tuyển sinh 10 TP.HCM năm học 2025-2026 – Bài 7
Từ một điểm $A$ nằm ngoài đường tròn $(O; R)$ với $OA = 2R$, kẻ hai tiếp tuyến $AB, AC$ đến đường tròn ($B, C$ là các tiếp điểm). Vẽ đường kính $BD$ của đường tròn $(O)$. Gọi $E$ là giao điểm thứ hai của đường thẳng $AD$ với đường tròn $(O)$. Đường thẳng $BC$ và $AO$ cắt nhau tại $H$.
\begin{enumerate}[a)]
	\item Chứng minh rằng tam giác $BED$ vuông và $ABHE$ là tứ giác nội tiếp.
	\item Chứng minh rằng $OD^2 = OH \cdot OA$ và $\widehat{HDO} = \widehat{HBE}$.
	\item Tính theo $R$ chu vi và diện tích của tam giác $DHE$.
\end{enumerate}
\loigiai{
\begin{enumerate}[a)]
	\item Điểm $E$ thuộc đường tròn $(O)$ đường kính $BD$ nên $\widehat{BED} = 90^\circ$ (góc nội tiếp chắn nửa đường tròn).\\
	Suy ra tam giác $BED$ vuông tại $E$, do đó $BE \perp AD$ tại $E$, hay $\widehat{AEB} = 90^\circ$.\\
	Vì $AB, AC$ là hai tiếp tuyến của $(O)$ cắt nhau tại $A$ nên $AB = AC$.\\
	Mặt khác $OB = OC = R$ nên $AO$ là đường trung trực của đoạn thẳng $BC$, suy ra $AO \perp BC$ tại $H$, hay $\widehat{AHB} = 90^\circ$.\\
	Xét tứ giác $ABHE$, ta có $\widehat{AEB} = 90^\circ$ và $\widehat{AHB} = 90^\circ$.\\
	Hai đỉnh $E$ và $H$ kề nhau cùng nhìn cạnh $AB$ dưới một góc vuông, do đó $ABHE$ là tứ giác nội tiếp đường tròn đường kính $AB$.
	\item Xét tam giác $ABO$ vuông tại $B$ có đường cao $BH$. Áp dụng hệ thức lượng trong tam giác vuông:
	$OB^2 = OH \cdot OA$.\\
	Mà $OB = OD = R$ nên $OD^2 = OH \cdot OA$.\\
	Từ $OD^2 = OH \cdot OA$, suy ra $\dfrac{OD}{OH} = \dfrac{OA}{OD}$.\\
	Xét tam giác $ODH$ và tam giác $OAD$, ta có:
	$\dfrac{OD}{OH} = \dfrac{OA}{OD}$ và góc $\widehat{AOD}$ chung.\\
	Suy ra $\Delta ODH \backsim \Delta OAD$ (c.g.c), do đó $\widehat{ODH} = \widehat{OAD}$, hay $\widehat{HDO} = \widehat{EAB}$.\\
	Mặt khác, trong tứ giác nội tiếp $ABHE$, ta có $\widehat{HBE} = \widehat{EAH}$ (hai góc nội tiếp cùng chắn cung $EH$).\\
	Mà $E, A, D$ thẳng hàng nên $\widehat{EAH} = \widehat{EAB}$.\\
	Từ đó suy ra $\widehat{HDO} = \widehat{HBE}$.
	\item Trong tam giác vuông $ABO$ tại $B$ có $OA = 2R$ và $OB = R$:\\
	$\cos\widehat{BOA} = \dfrac{OB}{OA} = \dfrac{R}{2R} = \dfrac{1}{2} \Rightarrow \widehat{BOA} = 60^\circ$.\\
	$AB = \sqrt{OA^2 - OB^2} = \sqrt{(2R)^2 - R^2} = R\sqrt{3}$.\\
	Vì $BD$ là đường kính nên $BD = 2R$ và $BD \perp AB$ tại $B$. Xét tam giác vuông $ABD$ tại $B$:
	$AD = \sqrt{AB^2 + BD^2} = \sqrt{(R\sqrt{3})^2 + (2R)^2} = \sqrt{3R^2 + 4R^2} = R\sqrt{7}$.\\
	Vì tam giác $ABD$ vuông tại $B$ có đường cao $BE$ nên:
	$BD^2 = DE \cdot DA \Rightarrow DE = \dfrac{BD^2}{AD} = \dfrac{4R^2}{R\sqrt{7}} = \dfrac{4\sqrt{7}}{7}R$.\\
	$BE = \sqrt{BD^2 - DE^2} = \sqrt{4R^2 - \dfrac{16}{7}R^2} = \dfrac{2\sqrt{21}}{7}R$.\\
	Trong tam giác vuông $OBH$ tại $H$ có $\widehat{BOH} = 60^\circ$:
	$BH = OB \cdot \sin 60^\circ = \dfrac{R\sqrt{3}}{2}$.\\
	Xét tam giác $HDB$ và tam giác $HBE$ có:
	$\widehat{HDB} = \widehat{HBE}$ (chứng minh ở câu b) và $\widehat{HEB} = \widehat{HBD}$ (cùng bằng góc $\widehat{BAO} = 30^\circ$).\\
	Suy ra $\Delta HDB \backsim \Delta HBE$ (g.g).\\
	Do đó: $\dfrac{BE}{BD} = \dfrac{HE}{HB} = \dfrac{HB}{HD}$.\\
	Thay số: $\dfrac{\dfrac{2\sqrt{21}}{7}R}{2R} = \dfrac{\sqrt{21}}{7}$.\\
	Suy ra:
	$HE = HB \cdot \dfrac{\sqrt{21}}{7} = \dfrac{R\sqrt{3}}{2} \cdot \dfrac{\sqrt{21}}{7} = \dfrac{3\sqrt{7}}{14}R$.\\
	$HD = \dfrac{HB}{\dfrac{\sqrt{21}}{7}} = \dfrac{\dfrac{R\sqrt{3}}{2}}{\dfrac{\sqrt{21}}{7}} = \dfrac{\sqrt{7}}{2}R$.\\
	Chu vi tam giác $DHE$ là:
	$P = DE + HE + HD = \dfrac{4\sqrt{7}}{7}R + \dfrac{3\sqrt{7}}{14}R + \dfrac{\sqrt{7}}{2}R = \left(\dfrac{8 + 3 + 7}{14}\right)\sqrt{7}R = \dfrac{18\sqrt{7}}{14}R = \dfrac{9\sqrt{7}}{7}R$.\\
	Trong tam giác $HED$, ta có $\widehat{HED} = 180^\circ - \widehat{AEB} - \widehat{BEH} = 90^\circ - \widehat{BEH}$. Mặt khác do tứ giác $ABHE$ nội tiếp nên $\widehat{BEH} = \widehat{BAH} = 30^\circ$, suy ra $\widehat{HED} = 60^\circ$.\\
	Diện tích tam giác $DHE$ là:
	$S = \dfrac{1}{2} \cdot HE \cdot DE \cdot \sin\widehat{HED} = \dfrac{1}{2} \cdot \dfrac{3\sqrt{7}}{14}R \cdot \dfrac{4\sqrt{7}}{7}R \cdot \sin 60^\circ = \dfrac{1}{2} \cdot \dfrac{12\cdot 7}{98}R^2 \cdot \dfrac{\sqrt{3}}{2} = \dfrac{3\sqrt{3}}{14}R^2$.\\
	Vậy chu vi tam giác $DHE$ là $\dfrac{9\sqrt{7}}{7}R$ và diện tích là $\dfrac{3\sqrt{3}}{14}R^2$.
\end{enumerate}
}
\end{bt}
